\begin{problem}[Two dice and a sum of at least $10$]
Two fair six-sided dice are rolled. Find the probability that the sum is at least $10$.
\end{problem}

\begin{solution}
There are $36$ equally likely ordered outcomes.

The sums at least $10$ are:
\[
10:(4,6),(5,5),(6,4),
\]
\[
11:(5,6),(6,5),
\]
\[
12:(6,6).
\]
So there are $6$ favourable outcomes.

Hence
\[
P(\text{sum at least }10)=\frac{6}{36}=\frac16.
\]
\end{solution}

\begin{takeaways}
\item Start by checking whether outcomes are equally likely.
\item With dice, ordered pairs are usually the cleanest sample space.
\item ``At least'' often means grouping several cases together.
\end{takeaways}

\begin{problem}[A card that is a heart or a face card]
A card is chosen at random from a standard deck of $52$ cards. Find the probability that the card is a heart or a face card.
\end{problem}

\begin{solution}
Let
\[
H=\{\text{heart}\},
\qquad
F=\{\text{face card}\}.
\]

There are $13$ hearts and $12$ face cards. The overlap consists of the heart face cards:
\[
J\heartsuit,\ Q\heartsuit,\ K\heartsuit,
\]
so
\[
n(H\cap F)=3.
\]

By the addition rule,
\[
P(H\cup F)=\frac{13+12-3}{52}=\frac{22}{52}=\frac{11}{26}.
\]
\end{solution}

\begin{takeaways}
\item Use the addition rule when two descriptions can overlap.
\item Subtract the intersection once to avoid double-counting.
\item Drawing a quick Venn diagram is often helpful.
\end{takeaways}

\begin{problem}[At least one blue marble]
A bag contains $4$ red marbles and $3$ blue marbles. Two marbles are drawn without replacement. Find the probability that at least one marble is blue.
\end{problem}

\begin{solution}
The complement of ``at least one blue'' is ``no blue'', which means both marbles are red.

So
\[
P(\text{both red})=\frac{4}{7}\cdot\frac{3}{6}=\frac{12}{42}=\frac27.
\]

Therefore,
\[
P(\text{at least one blue})=1-\frac27=\frac57.
\]
\end{solution}

\begin{takeaways}
\item ``At least one'' is often easiest through a complement.
\item Without replacement, multiply changing fractions.
\item The complement method reduces casework.
\end{takeaways}

\begin{problem}[Mutually exclusive or independent?]
One card is drawn from a standard deck. Let
\[
A=\{\text{the card is an ace}\},
\qquad
B=\{\text{the card is a spade}\}.
\]
Find $P(A\cap B)$ and decide whether $A$ and $B$ are independent.
\end{problem}

\begin{solution}
The only card that is both an ace and a spade is the ace of spades, so
\[
P(A\cap B)=\frac1{52}.
\]

Also,
\[
P(A)=\frac4{52}=\frac1{13},
\qquad
P(B)=\frac{13}{52}=\frac14.
\]

Thus
\[
P(A)P(B)=\frac1{13}\cdot\frac14=\frac1{52}.
\]
Since
\[
P(A\cap B)=P(A)P(B),
\]
the events are independent.
\end{solution}

\begin{takeaways}
\item Events can overlap and still be independent.
\item Independence is a probability test, not a visual one.
\item A single shared outcome does not automatically mean dependence.
\end{takeaways}
